\subsection{2. Hướng phát triển}
\begin{itemize}[label={--}, itemsep=1pt]
    \item \textbf{Mở rộng phạm vi thực nghiệm:} Tiếp tục thử nghiệm kiến trúc đề xuất trên các họ mạng trích xuất đặc trưng hiện đại khác (như EfficientNet, ConvNeXt hoặc Vision Transformer) để đánh giá khách quan hơn về tính ổn định và khả năng thích nghi của module MGAT + FastKAN/MLP đối với các dạng đặc trưng đầu vào khác nhau.
    
    \item \textbf{Làm giàu tập dữ liệu:} Tìm kiếm và thu thập thêm dữ liệu video văn hóa từ nhiều nguồn khác nhau để gia tăng kích thước tập huấn luyện. Việc này không chỉ giúp mô hình học được các đặc trưng tổng quát hơn mà còn đảm bảo kết quả đánh giá mang tính khách quan cao.
    
    \item \textbf{Nâng cao khả năng suy luận của đồ thị:} Nghiên cứu các phương pháp xây dựng đồ thị đa phương thức có tính giải thích cao hơn (Explainable AI). Thay vì chỉ học các mối quan hệ dựa trên độ tương đồng hoặc sự đối kháng thông qua trọng số cạnh, hướng phát triển tiếp theo sẽ tập trung vào việc mô hình hóa các quan hệ nhân quả hoặc ngữ nghĩa tường minh (Reasoning Graph), giúp máy tính thực sự "hiểu" được nội hàm văn hóa của di sản.
\end{itemize}